\documentclass[a4paper,11pt]{article}

% ── Encodage & langue ──────────────────────────────────────────────────────────
\usepackage[utf8]{inputenc}
\usepackage[T1]{fontenc}
\usepackage[french]{babel}

% ── Mise en page ───────────────────────────────────────────────────────────────
\usepackage[top=2.5cm, bottom=2.5cm, left=2.5cm, right=2.5cm]{geometry}
\usepackage{microtype}
\usepackage{parskip}

% ── Mathématiques ──────────────────────────────────────────────────────────────
\usepackage{amsmath, amssymb}
\usepackage{mathtools}

% ── Graphiques ─────────────────────────────────────────────────────────────────
\usepackage{graphicx}

% ── En-têtes & pieds de page ───────────────────────────────────────────────────
\usepackage{fancyhdr}
\usepackage{etoolbox}

% style des pages normales
\fancypagestyle{main}{%
  \fancyhf{}
  \renewcommand{\headrulewidth}{0.4pt}
  \renewcommand{\footrulewidth}{0.4pt}
  \lhead{\includegraphics[height=1cm]{logo_centrale.png}}
  \rhead{}
  \cfoot{\small Transition écologique \hfill \thepage}
}

% style page de titre (pas d'en-tête)
\fancypagestyle{titre}{%
  \fancyhf{}
  \renewcommand{\headrulewidth}{0pt}
  \renewcommand{\footrulewidth}{0pt}
}

\setlength{\headheight}{40pt}

% ── Couleurs ───────────────────────────────────────────────────────────────────
\usepackage{xcolor}
\definecolor{centralered}{RGB}{160,20,30}

% ── Polices & titres ───────────────────────────────────────────────────────────
% Titres de sections en small caps, style sobre
\usepackage{titlesec}
\titleformat{\section}[block]
  {\large\bfseries}
  {\Roman{section}\quad}
  {0pt}{}[\vspace{2pt}\titlerule]
\titleformat{\subsection}[block]
  {\normalsize\bfseries}
  {\thesubsection\quad}
  {0pt}{}

% ── Boîtes ─────────────────────────────────────────────────────────────────────
\usepackage{tcolorbox}
\tcbuselibrary{skins, breakable}

\newtcolorbox{remarque}[1][]{
  colback=gray!6!white, colframe=gray!50!black,
  fonttitle=\bfseries\small, title=Remarque,
  boxrule=0.4pt, #1
}

\newtcolorbox{piste}[1][]{
  colback=gray!4!white, colframe=centralered!70!black,
  fonttitle=\bfseries\small, title=Piste pour l'an prochain,
  boxrule=0.4pt, #1
}

\newtcolorbox{formule}[1][]{
  colback=gray!4!white, colframe=gray!60!black,
  fonttitle=\bfseries\small, title=Formule,
  boxrule=0.4pt, #1
}

% ── Hyperliens ─────────────────────────────────────────────────────────────────
\usepackage[hidelinks]{hyperref}

% ══════════════════════════════════════════════════════════════════════════════
\begin{document}

% ── Page de titre ──────────────────────────────────────────────────────────────
\thispagestyle{titre}

\begin{center}
  \includegraphics[width=7cm]{logo_centrale.png}\\[2em]
  {\large\textsc{CapECL}}\\[0.5em]
  {\large\textsc{Transition écologique}}\\[3em]
  \rule{\linewidth}{0.8pt}\\[1em]
  {\LARGE\textbf{Projet CREP 2025-2026}}\\[0.3em]
  {\large Annexe — Pistes et résultats utiles pour l'an prochain}\\[1em]
  \rule{\linewidth}{0.8pt}\\[3em]
\end{center}

\begin{minipage}[t]{0.5\textwidth}
  \textit{Élèves :}\\[0.4em]
  Camille Simond\\
  Maxence Lucas Morel\\
  Bilel Bouchama\\
  Elsa Mauron\\
  Lilian Laffont\\
  Mathéo Dhenaut\\
  Vittoria Gelot\\
  Anatole Ridereau\\
  Lucas Poirot\\
  Oïhana Daguerre\\
  Isshac Taibi
\end{minipage}%
\begin{minipage}[t]{0.5\textwidth}
  \raggedleft
  \textit{Enseignants :}\\[0.4em]
  M. Chevereau\\
  M. Bernard
\end{minipage}

\newpage

% ── Table des matières ─────────────────────────────────────────────────────────
\pagestyle{main}
\tableofcontents
\newpage

% ══════════════════════════════════════════════════════════════════════════════
\section{Présentation de l'annexe}

Ce document rassemble les résultats et pistes de travail jugés utiles à conserver
pour les prochaines séances, voire pour les groupes des années suivantes. Il est
tenu à part du carnet de bord afin de garder celui-ci concis, conformément aux
retours reçus.

On distingue deux catégories de contenu :
\begin{itemize}
  \item les formules et modèles \textbf{explorés mais finalement non utilisés},
        conservés car ils pourraient servir de point de départ à une modélisation
        plus ambitieuse ;
  \item les \textbf{pistes d'approfondissement} identifiées en cours de projet,
        destinées aux groupes futurs.
\end{itemize}

% ══════════════════════════════════════════════════════════════════════════════
\section{Formule de température en fonction de l'altitude — \texorpdfstring{$T(z)$}{T(z)}}

\subsection{Formule}

\begin{formule}
\begin{equation}
  T(z) = T_{\mathrm{eff}} \cdot \left[1 + \tau_0 \cdot e^{-\,z/H_{0,\mathrm{GES}}}\right]^{1/4}
  \label{eq:Tz_serre}
\end{equation}
\end{formule}

où :
\begin{itemize}
  \item $T_{\mathrm{eff}}$ : température effective de la Terre en l'absence
        d'atmosphère ($\approx 255\,\mathrm{K}$) ;
  \item $\tau_0$ : épaisseur optique totale des GES au niveau du sol,
        décomposable par gaz : $\tau_0 = \sum_i \tau_{0,i}$ ;
  \item $H_{0,\mathrm{GES}}$ : hauteur d'échelle des GES (décroissance
        exponentielle avec l'altitude) ;
  \item $z$ : altitude (en km ou en m selon l'unité choisie pour $H_{0,\mathrm{GES}}$).
\end{itemize}

% ══════════════════════════════════════════════════════════════════════════════
\section{Démonstration alternative de la formule de température}

\subsection{Pourquoi la garder}

Une démonstration de la formule de température a été établie à partir
de la loi de Stefan-Boltzmann et d'un modèle simplifié de l'effet de serre.
Cette formule reste valable pour la troposphère et présente plusieurs avantages :

\begin{itemize}
  \item \textbf{Modularité} : $\tau_0$ se décompose par GES, ce qui permet de
        paramétrer indépendamment la contribution de chaque gaz
        ($\mathrm{H_2O}$, $\mathrm{CO_2}$, $\mathrm{CH_4}$, $\mathrm{O_3}$\ldots).
  \item \textbf{Cohérence physique} : au sol ($z = 0$),
        $T(0) = T_{\mathrm{eff}} \cdot (1 + \tau_0)^{1/4} > T_{\mathrm{eff}}$ —
        l'effet de serre réchauffe bien la surface. En altitude,
        $T(z) \to T_{\mathrm{eff}}$ quand $\tau(z) \to 0$.
\end{itemize}

\subsection{Démonstration}

D'après la loi de Stefan-Boltzmann, le flux radiatif émis par une surface à la
température $T$ vaut :
\begin{equation}
  F = \sigma T^4
\end{equation}

Dans un modèle atmosphérique simplifié, l'absorption du rayonnement infrarouge
par les GES est représentée par un facteur correctif $(1 + \varepsilon)$, où
$\varepsilon$ caractérise l'intensité de l'effet de serre. Le flux effectivement
évacué vers l'espace s'écrit alors :
\begin{equation}
  F = \frac{\sigma\, T(z)^4}{1 + \varepsilon(z)}
\end{equation}

À l'équilibre radiatif, ce flux est égal au flux associé à la température effective :
\begin{equation}
  \sigma T_{\mathrm{eff}}^4 = \frac{\sigma\, T(z)^4}{1 + \varepsilon(z)}
\end{equation}

Après simplification par $\sigma$ :
\begin{equation}
  T(z)^4 = T_{\mathrm{eff}}^4 \cdot \bigl(1 + \varepsilon(z)\bigr)
  \implies
  T(z) = T_{\mathrm{eff}} \cdot \bigl(1 + \varepsilon(z)\bigr)^{1/4}
\end{equation}

En supposant que l'intensité de l'effet de serre décroît exponentiellement avec
l'altitude :
\begin{equation}
  \varepsilon(z) = \tau_0 \exp\!\left(-\frac{z}{H_{0,\mathrm{GES}}}\right)
\end{equation}

on retrouve la formule~\eqref{eq:Tz_serre}.

\subsection{À faire pour la réutiliser}

\begin{piste}
\begin{itemize}
  \item Préciser les valeurs de $\tau_{0,i}$ et $H_{0,\mathrm{GES}}$ pour chaque GES.
  \item Vérifier la cohérence des unités (altitude en m ou en km).
\end{itemize}
\end{piste}

\subsection{Lien entre la vapeur d'eau et l'effet de serre}

La vapeur d'eau est le principal gaz à effet de serre naturel de l'atmosphère.
Sa concentration étant maximale près de la surface, elle contribue fortement à
l'absorption du rayonnement infrarouge terrestre dans les basses couches de la
troposphère.

Comme la masse volumique de vapeur d'eau décroît exponentiellement avec
l'altitude, son influence sur l'effet de serre diminue également lorsque l'on
s'élève dans l'atmosphère. En première approximation :
\begin{equation}
  \tau_{\mathrm{H_2O}}(z)
  = \tau_{\mathrm{H_2O},0}\,\exp\!\left(-\frac{z}{H_{0,\mathrm{H_2O}}}\right)
\end{equation}

où $\tau_{\mathrm{H_2O},0}$ représente l'épaisseur optique de la vapeur d'eau
au niveau du sol. Cette relation présente la même forme que celle de la masse
volumique de vapeur d'eau, traduisant le fait que l'absorption du rayonnement
infrarouge dépend directement de la quantité de vapeur d'eau présente.

Ainsi, les premiers kilomètres de la troposphère concentrent l'essentiel de la
contribution de la vapeur d'eau à l'effet de serre terrestre.

% ══════════════════════════════════════════════════════════════════════════════
\section{Formules explorées et non utilisées au-delà de la troposphère}
\label{sec:abandonnees}

\subsection{Contexte}

Dans l'objectif de modéliser le profil de température $T(z)$ jusqu'à
$80\,\mathrm{km}$ d'altitude, nous avons cherché à étendre la formule issue
de l'équilibre radiatif-convectif~\eqref{eq:Tz_serre} au-delà de la
tropopause. Cette tentative a été abandonnée : dans la stratosphère, la
température \emph{remonte} avec l'altitude en raison de l'absorption des
rayonnements ultraviolets par l'ozone, un phénomène incompatible avec le
modèle convectif simplifié. Un modèle descriptif affine par morceaux
(type USS~1976) a finalement été retenu à la place.

Les deux formules ci-dessous sont conservées ici comme \textbf{points de départ
documentés} pour les groupes futurs qui souhaiteraient tenter une modélisation
physique explicite de la stratosphère.

\subsection{Formules}

\begin{formule}[title=Formule 1 — Profil de température avec effet de serre]
\begin{equation}
  T(z) = T_{\mathrm{eff}} \cdot \left[1 + \tau_0 \cdot e^{-\,z/H_{0,\mathrm{GES}}}\right]^{1/4}
\end{equation}
Valide uniquement dans la troposphère (jusqu'à $\approx 11\,\mathrm{km}$).
Non extensible à la stratosphère sans reformulation complète du terme d'absorption,
car la montée en température due à l'ozone n'y est pas captée.
\end{formule}

\begin{formule}[title=Formule 2 — Profil de masse volumique (loi barométrique)]
\begin{equation}
  \rho(z) = \rho_0\, e^{-z/h_0}
\end{equation}
où $\rho_0$ est la masse volumique au sol et $h_0 \approx 8\,\mathrm{km}$
l'échelle de hauteur de l'atmosphère.
Cette approximation exponentielle est satisfaisante dans la troposphère, mais
sous-estime la densité dans la stratosphère et au-delà, où la température n'est
plus uniforme.
\end{formule}

% ══════════════════════════════════════════════════════════════════════════════
\section{Pistes d'approfondissement pour les groupes futurs}

\begin{piste}[title=Modélisation physique de la stratosphère]
La remontée de température entre $20$ et $47\,\mathrm{km}$ est due à
l'absorption des rayonnements UV ($200$–$300\,\mathrm{nm}$) par l'ozone
($\mathrm{O_3}$). Un modèle plus ambitieux pourrait :
\begin{enumerate}
  \item Calculer l'épaisseur optique de l'ozone $\tau_{\mathrm{O_3}}(z)$ en
        fonction de l'altitude, à partir du profil de concentration en ozone.
  \item Utiliser la loi de Beer-Lambert pour estimer le flux UV absorbé à
        chaque altitude.
  \item En déduire un terme de chauffage $Q(z)$ et résoudre l'équilibre
        radiatif dans la stratosphère.
\end{enumerate}
Cette approche permettrait d'obtenir un modèle \emph{explicatif} (et non plus
seulement descriptif) sur toute la colonne atmosphérique.
\end{piste}

\begin{piste}[title=Couplage avec les GES variables dans le temps]
La formule~\eqref{eq:Tz_serre} est statique. Une extension naturelle serait
de faire varier les paramètres $\tau_{0,i}$ en fonction de l'année, en utilisant
les données historiques et les projections des concentrations en GES
($\mathrm{CO_2}$, $\mathrm{CH_4}$, $\mathrm{H_2O}$\ldots). Cela permettrait
de simuler l'évolution du profil vertical de température au cours du temps.
\end{piste}

\begin{piste}[title=Loi barométrique améliorée]
La loi $\rho(z) = \rho_0\, e^{-z/h_0}$ pourrait être utilisée dans la
troposphère conjointement avec le profil de température pour calculer la
pression. Pour aller plus loin, le modèle standard retenu fournit des valeurs
de pression et de masse volumique par couche, qui pourraient être implémentées
selon le même principe affine par morceaux que $T(z)$.
\end{piste}

% ══════════════════════════════════════════════════════════════════════════════
\section*{Références}

\begin{itemize}
  \item NOAA/NASA/USAF, \emph{U.S. Standard Atmosphere 1976}, octobre 1976.
        \href{https://www.ngdc.noaa.gov/stp/space-weather/online-publications/miscellaneous/us-standard-atmosphere-1976/us-standard-atmosphere_st76-1562_noaa.pdf}{\underline{ici}}
  \item M.~Cavcar, \emph{The International Standard Atmosphere (ISA)},
        Anadolu University.
        \href{http://fisicaatmo.at.fcen.uba.ar/practicas/ISAweb.pdf}{\underline{ici}}

\end{itemize}

\end{document}